\uuid{o2Ta}
\titre{Équation droite/plan}
\auteur{bodin ; hanani ; mzari}
\organisation{exo7}
\licence{CC BY-NC-SA 4.0}
\datecreate{2020-01-20}
\niveau{L1}
\module{Géométrie}
\chapitre{Géométrie de l'espace}
\sousChapitre{Plans}
\difficulte{2}
\theme{Plans -- Droites, équation droite, plan}
\exo7qcmid{lille-1.geometrie-espace.012}

\contenu{
\texte{
Pour ces questions, l'espace est muni d'un repère orthonormé direct  $(O, \vec{i}, \vec{j}, \vec{k})$.
}

\begin{question}
\qclassification{141.03,141.04}
\qtags{motcle=équation droite/plan}
Soit $P : x+y-z=1$,  $P' : x+z=-1$ deux plans et $Q$ le plan passant par $A(1,1,1)$ et perpendiculaire à $P$ et à $P'$. Quelles sont les assertions vraies ?
\begin{answers}

    \bad{Une équation cartésienne de $Q$ est $x+2y-z+2=0$.}

    \good{Une équation cartésienne de $Q$ est $x-2y-z+2=0$.}

   \bad{Une représentation paramétrique de $Q$ est :
 $$\left\{\begin{array}{ccl}x&=&1-t\\y&=&1+t-s\\ z&=&1+t+s,\quad (t,s \in \Rr) \end{array}\right.$$}

    \good{Une représentation paramétrique de $Q$ est :
 $$\left\{\begin{array}{ccl}x&=&1+t+s\\y&=&1+t\\ z&=&1-t+s,\quad (t,s \in \Rr) \end{array}\right.$$}

\end{answers}

\begin{explanations}
$Q$ passe par $A$ et est dirigé par $\vec{n}$ et $\vec{n'}$, où $\vec{n}$ et $\vec{n'}$ sont des vecteurs normaux à $P$ et à $P'$ respectivement.
\end{explanations}
\end{question}
}
